\section{LLM 为何如此令人兴奋}

\subsection{参数规模的爆炸式增长}

如果 n-gram 语言模型早在 1980 年代就已存在，为什么人们直到近年才对语言模型如此兴奋？关键变化之一是\textbf{参数规模}的指数级增长：

\begin{table}[H]
\centering
\begin{tabular}{lll}
\toprule
\textbf{模型} & \textbf{年份} & \textbf{参数量} \\
\midrule
BERT Large     & 2018 & 3.4 亿（340M）  \\
GPT-3          & 2020 & 1750 亿（175B） \\
当代前沿模型   & 2024--2025 & 万亿级（trillions） \\
\bottomrule
\end{tabular}
\caption{语言模型参数规模的增长趋势}
\end{table}

参数可以理解为神经网络中的权重，类比于大脑中突触之间的连接。更多的参数意味着模型有更大的容量来``记忆''和``理解''关于世界的信息。

\subsection{上下文窗口的扩展}

另一个关键变化是\textbf{上下文窗口}（context window）的剧烈增长：

\begin{table}[H]
\centering
\begin{tabular}{lll}
\toprule
\textbf{模型架构} & \textbf{典型上下文长度} & \textbf{量级} \\
\midrule
N-gram 模型       & $\sim$4 个词          & 个位数 \\
基础 RNN           & $\sim$20 个词         & 十位数 \\
LSTM               & $\sim$200 个词        & 百位数 \\
早期 Transformer   & $\sim$2,048 tokens    & 千位数 \\
Gemini（2024+）    & $\sim$2,000,000 tokens & 百万级 \\
\bottomrule
\end{tabular}
\caption{上下文窗口的历史演变}
\end{table}

前面 n-gram 模型的概率循环问题，正是因为上下文窗口只有 4 个词。当窗口扩大到百万 token 量级，模型能够``看到''和处理的信息量产生了质变。

\subsection{Few-Shot Learning 的涌现}

\begin{importantbox}{GPT-3 论文的核心发现}
2020 年的 GPT-3 论文（``Language Models are Few-Shot Learners''）报告了一个关键的涌现行为（emergent behavior）：当模型参数量达到约 1750 亿时，模型开始展现出成功的 zero-shot、one-shot 和 few-shot prompting 能力——无需任何梯度更新（no gradient updates），仅通过 prompt 中的指令或少量示例，就能泛化到全新的任务。
\end{importantbox}

这三种 prompting 方式的定义：

\begin{itemize}
    \item \textbf{Zero-shot}：只给指令，不给示例。例如：``Translate English to French. Cheese:''
    \item \textbf{One-shot}：给一个示例。例如：``Sea otter $\rightarrow$ Loutre de mer. Cheese $\rightarrow$''
    \item \textbf{Few-shot}：给若干示例（通常 3--5 个）
\end{itemize}

GPT-3 论文的核心图表显示：在低参数量时，zero/one/few-shot 性能接近随机；当参数量超过某个临界点后，性能急剧上升。这意味着大规模语言模型获得了一种类似人类``举一反三''的能力——给出少量示例就能泛化到新情况。

\begin{knowledgebox}{``No Gradient Updates'' 的意义}
传统方法要让模型掌握新任务，需要通过 fine-tuning 更新权重（gradient updates）。GPT-3 论文特别强调了 ``no gradient updates''，因为 few-shot learning 完全在推理时（inference time）完成，无需修改模型参数。这大幅降低了将 LLM 应用到新任务的成本——你只需要写好 prompt。
\end{knowledgebox}

讲师同时指出了 Chinchilla 论文的贡献：通过更高效的训练策略（数据与计算的最优分配），可以用更少的参数达到相似的性能。这带来了更低的功耗、更快的推理和更小的内存占用。但研究者很快将节省下来的资源用于继续 scale up，推动参数量从千亿级进入万亿级。

\subsection{本章小结}

LLM 令人兴奋的根本原因是三大变化的叠加：(1) 参数量从百万级增长到万亿级；(2) 上下文窗口从个位数增长到百万级；(3) 在足够大的规模下涌现出了无需微调即可泛化的 few-shot learning 能力。这三者共同推动了从``笨拙的自动补全''到``通用任务求解器''的质变。

